\section{Conclusion and Political Economy Synthesis}
\label{sec:conclusion}

\textbf{6.1} For over half a century, the economic history of the Chilean Unidad Popular has been framed as an open-and-shut case of macroeconomic populism and unforced policy failure. By treating the Central Bank as an unconstrained printing press and inflation as an unmediated consequence of fiscal deficits, the orthodox monetarist and Austrian traditions \citep{Edwards2024,EspinosaVejar2025} constructed a powerful narrative of self-inflicted collapse. As our epistemological critique showed, this orthodoxy rests on an illusion of causality \citep{BlancoMatute2018}, narrative license \citep{IschEtAl2026}, and an M-bias collider trap that mistook the symptoms of private distribution lockouts and external blockades for the primary drivers of crisis.

\textbf{6.2} By formulating the empirical inquiry as a formal Neo-Lakatosian Hypothesis Tournament \citep{KingSznajder2006}, this chapter submitted the foundational claims of these competing schools to econometric adjudication:
\begin{enumerate}[leftmargin=2em]
	\item \textbf{Falsification of the Monetarist Tradition ($H_1$):} The core monetarist premise of exogenous money printing leading inflation ($g_M \to \pi$) is refuted across all empirical specifications. Instead, we find robust evidence of reverse accommodation ($\pi_{t-1} \to g_{M,t}, p < 0.001$) and an explicit accommodation gradient ($\chi^2_{M0} > \chi^2_{M1} > \chi^2_{M2}$), confirming that monetary expansion was an institutional Lender of Last Resort response to structural bottlenecks and foreign exchange strangulation.
	\item \textbf{Falsification of the Austrian Business-Cycle Tradition ($H_2$):} The Austrian claim that the 1971 boom reflected an unsustainable capital lengthening and roundabout malinvestment boom ($\Delta K \gg 0$) is refuted on its own empirical terrain: Gross Fixed Capital Formation contracted ($-2.3\%$) while capacity utilization hit an all-time peak of $98.44\%$, validating Keynesian-Structuralist idle capacity reactivation.
	\item \textbf{Validation of the Structural-Dependency Tradition ($H_3$):} High-frequency monthly TVAR estimates based on \citet{ClavijoCortesKappesRochon2025} confirm that monetary pass-through is strictly state-dependent. In Foley's (2003) Hedge regime ($\text{SolvR}_m > 0.62$), monetary expansion is non-inflationary ($\beta \approx 0.00\%$). It is only when external blockades, terms-of-trade shifts, and credit embargos drain foreign exchange reserves below the critical solvency threshold ($\text{SolvR}_m \le 0.28$) that the economy enters a Ponzi regime, where unbacked liquidity passes directly into hyperinflation (+4.54\% at $h=0$).
\end{enumerate}

\textbf{6.3} The collapse of the Unidad Popular was the peripheral expression of a global crisis of capitalism in the early 1970s. When the Nixon administration closed the gold window on August 15, 1971, it dismantled the fixed-rate anchor of Bretton Woods and exported U.S. stagflation to the periphery. In Chile, this global monetary shock interacted with the domestic structural contradictions of peripheral underdevelopment and deliberate imperial destabilization. The UP did not collapse because it pursued income redistribution; it collapsed because it attempted structural transformation without command over the world money required to settle international obligations.

\textbf{6.4} For contemporary peripheral economies, this historical reassessment offers profound lessons. Monetary sovereignty cannot be achieved through domestic accounting alone. In a hierarchical international monetary order, the external balance sheet—summarized by the Solvency Ratio—remains the fundamental structural constraint on sovereign policy space. Sustainable developmental transitions in the global South require strategic management of the external corridor, macroprudential capital controls, and coordinated efforts to dismantle structural subordination to core reserve currencies.